\documentclass{beamer}
\usetheme{Warsaw}

\usepackage[utf8]{inputenc}
\usepackage[T1]{fontenc}
\usepackage[ngerman]{babel}
\usepackage{fancybox}
\usepackage{multimedia}
\usepackage{subfig}
\usepackage{amsmath}
\usepackage{amssymb}
\usepackage{booktabs}
\usepackage{listings}
\usepackage{xcolor}
\usepackage{hyperref}
\hypersetup{
    colorlinks=true,
    urlcolor=blue
}

\lstset{
    basicstyle=\ttfamily\footnotesize,
    keywordstyle=\color{blue}\bfseries,
    commentstyle=\color{gray}\itshape,
    stringstyle=\color{red!70!black},
    showstringspaces=false,
    language=Python,
    frame=single,
    breaklines=true
}

\setbeamertemplate{footline}[frame number]

\begin{document}

\title[Stochastik]
{Wahrscheinlichkeitstheorie und Statistik (Stochastik)
\\[1ex]
Lineare Kongruenz-Generatoren}
\author[Dr. Johannes Riesterer]
{Dr.\ rer.\ nat.\ Johannes Riesterer}
\date{}
\subject{Stochastik}

\frame{\titlepage}


\begin{frame}
    \frametitle{Pseudozufallszahlen}
    \framesubtitle{Motivation}

\begin{block}{Problem}
Ein Computer ist deterministisch -- echte Zufallszahlen kann er nicht
erzeugen. Für Simulationen (Monte Carlo, Bootstrapping, Kryptographie,
Spiele, \ldots) benötigt man jedoch Folgen, die sich \emph{wie} Zufall
verhalten.
\end{block}

\begin{block}{Idee: Pseudozufallszahlen}
Eine deterministische Folge $X_0, X_1, X_2, \ldots$, die statistische
Tests auf Zufälligkeit besteht. Vorteile:
\begin{itemize}
    \item reproduzierbar (Seed),
    \item schnell berechenbar,
    \item geringer Speicherbedarf.
\end{itemize}
\end{block}

\end{frame}


\begin{frame}
    \frametitle{Linearer Kongruenz-Generator (LCG)}
    \framesubtitle{Definition}

\begin{block}{Rekursionsvorschrift}
Wähle
\begin{itemize}
    \item Modulus $m \in \mathbb{N}$, $m > 0$,
    \item Multiplikator $a \in \{1, \ldots, m-1\}$,
    \item Inkrement $c \in \{0, \ldots, m-1\}$,
    \item Startwert (Seed) $X_0 \in \{0, \ldots, m-1\}$.
\end{itemize}
Die Folge ist definiert durch
\[
    X_{n+1} \;=\; (a \cdot X_n + c) \bmod m.
\]
\end{block}

\begin{block}{Normierung auf $[0,1)$}
\[
    U_n \;=\; \frac{X_n}{m} \;\in\; [0,1).
\]
\end{block}

\end{frame}


\begin{frame}
    \frametitle{Beispiel}
    \framesubtitle{$m=9,\ a=7,\ c=4,\ X_0=3$}

\begin{block}{Berechnung von Hand}
\[
\begin{array}{c|l|c}
n & \text{Rechnung} & X_n \\ \hline
0 & -                                       & 3 \\
1 & (7\cdot 3 + 4) \bmod 9 = 25 \bmod 9     & 7 \\
2 & (7\cdot 7 + 4) \bmod 9 = 53 \bmod 9     & 8 \\
3 & (7\cdot 8 + 4) \bmod 9 = 60 \bmod 9     & 6 \\
4 & (7\cdot 6 + 4) \bmod 9 = 46 \bmod 9     & 1 \\
5 & (7\cdot 1 + 4) \bmod 9 = 11 \bmod 9     & 2 \\
6 & (7\cdot 2 + 4) \bmod 9 = 18 \bmod 9     & 0 \\
7 & (7\cdot 0 + 4) \bmod 9 =  4 \bmod 9     & 4 \\
8 & (7\cdot 4 + 4) \bmod 9 = 32 \bmod 9     & 5 \\
9 & (7\cdot 5 + 4) \bmod 9 = 39 \bmod 9     & 3 \\
\end{array}
\]
\end{block}

Nach $9 = m$ Schritten wiederholt sich die Folge -- maximale Periode erreicht.

\end{frame}


\begin{frame}
    \frametitle{Periode}
    \framesubtitle{Hull--Dobell-Bedingungen}

\begin{block}{Beobachtung}
Da $X_n \in \{0, 1, \ldots, m-1\}$, gibt es höchstens $m$ verschiedene
Werte. Spätestens nach $m$ Schritten wiederholt sich die Folge.
\end{block}

\begin{block}{Satz (Hull \& Dobell, 1962)}
Ein LCG hat genau dann die maximale Periode $m$, wenn:
\begin{enumerate}
    \item $\gcd(c, m) = 1$,
    \item für jeden Primfaktor $p$ von $m$ gilt $p \mid (a-1)$,
    \item ist $4 \mid m$, so auch $4 \mid (a-1)$.
\end{enumerate}
\end{block}

\begin{block}{Beispiele aus der Praxis}
\begin{tabular}{lrrr}
\toprule
                    & $m$       & $a$         & $c$ \\
\midrule
glibc \texttt{rand()}        & $2^{31}$  & 1103515245  & 12345 \\
Numerical Recipes   & $2^{32}$  & 1664525     & 1013904223 \\
\bottomrule
\end{tabular}
\end{block}

\end{frame}


\begin{frame}[fragile]
    \frametitle{Implementierung in Python}

\begin{lstlisting}
def lcg(seed, a=1664525, c=1013904223, m=2**32, n=10):
    x = seed
    for _ in range(n):
        x = (a * x + c) % m
        yield x / m   # in [0, 1)

print(list(lcg(seed=42, n=5)))
\end{lstlisting}

\begin{block}{Bemerkung}
Mit demselben Seed liefert der Generator stets dieselbe Folge --
unverzichtbar für reproduzierbare Simulationen.
\end{block}

\end{frame}


\begin{frame}
    \frametitle{Eigenschaften und Grenzen}

\begin{block}{Vorteile}
\begin{itemize}
    \item sehr schnell, eine Multiplikation und eine Addition pro Schritt,
    \item minimaler Speicherbedarf (ein Integer),
    \item deterministisch reproduzierbar.
\end{itemize}
\end{block}

\begin{block}{Nachteile}
\begin{itemize}
    \item \textbf{Marsaglia-Effekt:} Die Tupel
          $(X_n, X_{n+1}, \ldots, X_{n+k-1})$ liegen in $[0,1)^k$ auf
          wenigen Hyperebenen $\Rightarrow$ schlechte Verteilung in
          hohen Dimensionen.
    \item niederwertige Bits haben oft sehr kurze Perioden,
    \item \textbf{nicht kryptographisch sicher}: aus wenigen Ausgaben
          lassen sich $a, c, m$ rekonstruieren.
\end{itemize}
\end{block}

\begin{block}{Heutige Standards}
Mersenne Twister (\texttt{numpy.random}, Python \texttt{random}),
PCG, xorshift, \ldots
\end{block}

\end{frame}


\end{document}
